
\documentclass{article}
\usepackage{colortbl}
\usepackage{makecell}
\usepackage{multirow}
\usepackage{supertabular}

\begin{document}

\newcounter{utterance}

\centering \large Interaction Transcript for game `cladder', experiment `full\_v1.5\_default', episode 6916 with refineonly{-}Qwen3.5{-}4B.
\vspace{24pt}

{ \footnotesize  \setcounter{utterance}{1}
\setlength{\tabcolsep}{0pt}
\begin{supertabular}{c@{$\;$}|p{.15\linewidth}@{}p{.15\linewidth}p{.15\linewidth}p{.15\linewidth}p{.15\linewidth}p{.15\linewidth}}
    \# & \multicolumn{2}{c}{Player} && \multicolumn{2}{c}{Game Master} \\
    \hline

    \theutterance \stepcounter{utterance}  
    & & & \multicolumn{4}{p{0.6\linewidth}}{
        \cellcolor[rgb]{0.9,0.9,0.9}{
            \makecell[{{p{\linewidth}}}]{
                \texttt{\tiny{[P1$\langle$GM]}}
                \texttt{You are an expert in causal inference. The following question is not a typical commonsense query, but rather a meticulously designed question created by a professor specializing in causal inference, intended to assess the students' mastery of the course content.} \\
\texttt{Imagine a self{-}contained, hypothetical world with only the following conditions, and without any unmentioned factors or causal relationships: Jyka has a direct effect on hwax. Lirg has a direct effect on hwax. Method 1: We look at how jyka correlates with lirg case by case according to hwax. Method 2: We look directly at how jyka correlates with lirg in general. To understand how jyka affects lirg, is it more correct to use the Method 1 than Method 2?} \\
\texttt{Start your answer with 'Yes' or 'No', followed by additional reasoning or evidence to support your explanation.} \\
            }
        }
    }
     \\ \\

    \theutterance \stepcounter{utterance}  
    & & \multicolumn{4}{p{0.6\linewidth}}{
        \cellcolor[rgb]{0.95,0.95,0.95}{
            \makecell[{{p{\linewidth}}}]{
                \texttt{\tiny{[P1$\leftrightarrow$P]}}
                \texttt{{[}DRAFT{]}} \\
\texttt{No} \\
\\ 
\texttt{**Reasoning:**} \\
\\ 
\texttt{In this hypothetical world, we are given two specific causal relationships: \$Jyka \textbackslash{}to Hwax\$ and \$Lirg \textbackslash{}to Hwax\$. There is no mention of a direct causal link between \$Jyka\$ and \$Lirg\$. In causal inference, when two variables share a common effect (here, \$Hwax\$), they form a **collider structure** (\$Jyka \textbackslash{}to Hwax \textbackslash{}leftarrow Lirg\$).} \\
\\ 
\texttt{1.  **Analysis of Method 2 (Direct Correlation):** If we simply look at the correlation between \$Jyka\$ and \$Lirg\$ in general (without conditioning on \$Hwax\$), we are observing the marginal relationship. In a collider structure, variables are typically **independent** unless there is some other unmentioned common cause. Therefore, Method 2 would likely show no correlation (or a correlation of zero) between \$Jyka\$ and \$Lirg\$. Since there is no direct causal arrow between them, concluding that \$Jyka\$ affects \$Lirg\$ based on this method is incorrect; they are likely causally unrelated.} \\
\\ 
\texttt{2.  **Analysis of Method 1 (Correlation via Hwax):** Method 1 involves looking at the relationship between \$Jyka\$ and \$Lirg\$ *case by case according to \$Hwax\$*. This is equivalent to **conditioning on the collider** (\$Hwax\$). A fundamental principle of causal inference is that conditioning on a collider (or its descendants) induces a spurious association between its parents. This phenomenon is known as **collider bias** or **M{-}bias** (in a specific context). By fixing \$Hwax\$ at specific values, we restrict the flow of probability, creating a statistical dependence between \$Jyka\$ and \$Lirg\$ even though they have no causal link.} \\
\\ 
\texttt{3.  **Conclusion on "Affecting":** The question asks how to understand how \$Jyka\$ *affects* \$Lirg\$.} \\
\texttt{    *   Method 2 correctly reflects the absence of a direct causal effect (showing independence).} \\
\texttt{    *   Method 1 incorrectly creates a spurious correlation due to the collider bias.} \\
\\ 
\texttt{Therefore, it is **not** more correct to use Method 1. In fact, Method 1 is the exact opposite of what is needed to determine causal relationships in this structure; it actively obscures the truth by introducing a false dependency. To understand the causal effect (which is zero), one must avoid conditioning on the collider (\$Hwax\$), making Method 2 the logically superior approach for assessing the direct causal link (or lack thereof). Using Method 1 would lead to the erroneous conclusion that \$Jyka\$ and \$Lirg\$ are related.} \\
            }
        }
    }
    & \\ \\

    \theutterance \stepcounter{utterance}  
    & \multicolumn{4}{p{0.6\linewidth}}{
        \cellcolor[rgb]{0.561,0.737,0.561}{
            \makecell[{{p{\linewidth}}}]{
                \texttt{\tiny{[P1$\rangle$GM]}}
                \texttt{No} \\
\\ 
\texttt{**Reasoning:**} \\
\\ 
\texttt{In this hypothetical world, we are given two specific causal relationships: \$Jyka \textbackslash{}to Hwax\$ and \$Lirg \textbackslash{}to Hwax\$. There is no mention of a direct causal link between \$Jyka\$ and \$Lirg\$. In causal inference, when two variables share a common effect (here, \$Hwax\$), they form a **collider structure** (\$Jyka \textbackslash{}to Hwax \textbackslash{}leftarrow Lirg\$).} \\
\\ 
\texttt{1.  **Analysis of Method 2 (Direct Correlation):** If we simply look at the correlation between \$Jyka\$ and \$Lirg\$ in general (without conditioning on \$Hwax\$), we are observing the marginal relationship. In a collider structure, variables are typically **independent** unless there is some other unmentioned common cause. Therefore, Method 2 would likely show no correlation (or a correlation of zero) between \$Jyka\$ and \$Lirg\$. Since there is no direct causal arrow between them, concluding that \$Jyka\$ affects \$Lirg\$ based on this method is incorrect; they are likely causally unrelated.} \\
\\ 
\texttt{2.  **Analysis of Method 1 (Correlation via Hwax):** Method 1 involves looking at the relationship between \$Jyka\$ and \$Lirg\$ *case by case according to \$Hwax\$*. This is equivalent to **conditioning on the collider** (\$Hwax\$). A fundamental principle of causal inference is that conditioning on a collider (or its descendants) induces a spurious association between its parents. This phenomenon is known as **collider bias**. By fixing \$Hwax\$ at specific values, we restrict the flow of probability, creating a statistical dependence between \$Jyka\$ and \$Lirg\$ even though they have no causal link.} \\
\\ 
\texttt{3.  **Conclusion on "Affecting":** The question asks how to understand how \$Jyka\$ *affects* \$Lirg\$.} \\
\texttt{    *   Method 2 correctly reflects the absence of a direct causal effect (showing independence).} \\
\texttt{    *   Method 1 incorrectly creates a spurious correlation due to the collider bias.} \\
\\ 
\texttt{Therefore, it is **not** more correct to use Method 1. In fact, Method 1 is the exact opposite of what is needed to determine causal relationships in this structure; it actively obscures the truth by introducing a false dependency. To understand the causal effect (which is zero), one must avoid conditioning on the collider (\$Hwax\$), making Method 2 the logically superior approach for assessing the direct causal link (or lack thereof). Using Method 1 would lead to the erroneous conclusion that \$Jyka\$ and \$Lirg\$ are related.} \\
            }
        }
    }
    & & \\ \\

    \theutterance \stepcounter{utterance}  
    & & & \multicolumn{2}{p{0.3\linewidth}}{
        \cellcolor[rgb]{0.9,0.9,0.9}{
            \makecell[{{p{\linewidth}}}]{
                \texttt{\tiny{[GM$|$GM]}}
                \texttt{No} \\
            }
        }
    }
    & & \\ \\

    \theutterance \stepcounter{utterance}  
    & & & \multicolumn{2}{p{0.3\linewidth}}{
        \cellcolor[rgb]{0.9,0.9,0.9}{
            \makecell[{{p{\linewidth}}}]{
                \texttt{\tiny{[GM$|$GM]}}
                \texttt{no} \\
            }
        }
    }
    & & \\ \\

    \theutterance \stepcounter{utterance}  
    & & & \multicolumn{2}{p{0.3\linewidth}}{
        \cellcolor[rgb]{0.9,0.9,0.9}{
            \makecell[{{p{\linewidth}}}]{
                \texttt{\tiny{[GM$|$GM]}}
                \texttt{game\_result = WIN} \\
            }
        }
    }
    & & \\ \\

\end{supertabular}
}

\end{document}
